\chapter{Case Study Baseline (As-Is System)}
\label{chap:baseline}

This chapter describes the initial state of the Outsight Cloud platform before the migration design was formalized. The goal is to establish the baseline against which migration decisions are evaluated.

\section{Current Architecture and Data Flow}
The as-is platform handled LiDAR-related assets and simulations through a web application and backend API stack:
\begin{itemize}
    \item Frontend based on Vue 3.
    \item Backend based on Node.js/Express.
    \item Redis-backed session mechanisms.
    \item Containerized deployment with reverse-proxy routing.
    \item File storage through S3-compatible object storage services.
    \item Core business data strongly coupled to Airtable.
\end{itemize}

This architecture enabled fast delivery in earlier stages of product evolution, but it created a concentration of dependencies in external data services and manually maintained infrastructure conventions.

\section{Operational and Cost Baseline}
The baseline analysis highlighted two classes of constraints.

\subsection{Operational Constraints}
Initial platform analysis identified:
\begin{itemize}
    \item Limited traceability in deployment and infrastructure changes.
    \item Residual complexity in API routing conventions.
    \item Difficulty in isolating and validating behavior when replacing data backends.
    \item Dependency on non-uniform operational knowledge across team members.
\end{itemize}

\subsection{Economic Constraints}
From the baseline reports, recurring cost and dependency concentration on Airtable emerged as significant migration drivers. Airtable was not only a persistence layer but also part of operational workflows, which increased both direct cost exposure and long-term coupling risk.

\section{Technical Debt and Migration Drivers}
The migration need was not framed as a generic technology update, but as a response to specific debt patterns:
\begin{itemize}
    \item \textbf{Data-layer coupling:} application behavior strongly tied to Airtable schema and workflows.
    \item \textbf{Identity evolution pressure:} growing need for standards-based authentication and cleaner authorization boundaries.
    \item \textbf{Infrastructure reproducibility gaps:} insufficiently formalized deployment model for multi-environment evolution.
    \item \textbf{Validation limitations:} absence of a robust incremental strategy to verify parity during backend transition.
\end{itemize}

These drivers motivated a migration program with two design priorities: incremental risk reduction and explicit decision traceability.

